\chapter{Computational Number Theory and Cryptography}\label{cryptography}
\chapterstatus{draft}

Some arithmetic problems are easy to compute and others, as far as anyone knows, are hard. Testing whether a number with a thousand digits is prime takes a fraction of a second; factoring a product of
two primes of five hundred digits each is beyond all present computers. Raising to a power modulo $p$ is fast; recovering the exponent, the \emph{discrete logarithm}, seems to be slow. Public-key
cryptography is built on this asymmetry. This chapter measures the cost of the basic algorithms (Euclid, exponentiation, multiplication), proves the correctness and error bounds of the primality tests
of Fermat, Miller--Rabin and Lucas, presents the classical factoring and discrete logarithm algorithms, and describes the Diffie--Hellman, RSA and ElGamal systems with some attacks that show which
hypotheses matter. The RSA system itself was introduced in Chapter~\ref{number-theory} (Proposition~\ref{number-theory:proposition-rsa}).

\noindent\textbf{Prerequisites.} Chapter~\ref{number-theory} (Elementary Number Theory), Chapter~\ref{groups} (Groups), Chapter~\ref{fields} (Fields and Galois Theory),
Chapter~\ref{computability-and-complexity} (Computability and Complexity), Chapter~\ref{prime-numbers} (Prime Numbers) and Chapter~\ref{geometry-of-numbers} (Geometry of Numbers).

\section{The cost of arithmetic}\label{cryptography:section-algorithms}

\begin{notation}\label{cryptography:notation-cost}
The \emph{size} of a positive integer $n$ is its number of binary digits, about $\log_2n$. An algorithm on integers of size at most $k$ is \emph{polynomial time} if it uses $O(k^c)$ bit operations for a constant
$c$. We write $\log$ for the natural logarithm, $\gcd$ for the greatest common divisor and $U_n = (\ZZ/n\ZZ)^\times$ for the group of units modulo $n$, of order $\varphi(n)$; congruences and products modulo $n$ are
written in the usual multiplicative notation.
\end{notation}

Schoolbook addition of two $k$-bit numbers costs $O(k)$ bit operations and schoolbook multiplication or division with remainder $O(k^2)$.

\begin{theorem}[Lamé]\label{cryptography:theorem-lame}
If the Euclidean algorithm applied to $a > b \geq 1$ uses $k$ divisions, then $b \geq F_{k+1}$ and $a \geq F_{k+2}$, where $F_m$ are the Fibonacci numbers. Consequently $k \leq \log_\varphi b + 1$ with
$\varphi = \frac{1 + \sqrt5}{2}$, and $\gcd(a, b)$, together with $u$, $v$ with $au + bv = \gcd(a, b)$, is computed in polynomial time.
\end{theorem}

\begin{proof}
Write $r_0 = a$, $r_1 = b$ and $r_{i-1} = q_ir_i + r_{i+1}$ for $i = 1, \dots, k$, with $r_1 > r_2 > \dots > r_k > r_{k+1} = 0$ and $q_i \geq 1$. Then $r_k \geq 1 = F_2$, $r_{k-1} > r_k$ gives $r_{k-1} \geq 2 = F_3$,
and $r_{i-1} \geq r_i + r_{i+1}$ gives by descending induction $r_i \geq F_{k+2-i}$; in particular $b = r_1 \geq F_{k+1}$ and $a \geq F_{k+2}$. By induction $F_m \geq \varphi^{m-2}$ (as $\varphi^2 = \varphi + 1$), so
$b \geq \varphi^{k-1}$. Each step costs $O(\log^2a)$, and keeping track of the coefficients (extended Euclid) multiplies the cost by a constant.
\end{proof}

\begin{example}\label{cryptography:example-lame}
Consecutive Fibonacci numbers are the worst case: $\gcd(F_{k+2}, F_{k+1})$ takes exactly $k$ steps, all quotients being $1$, for instance $(13, 8) \to (8, 5) \to (5, 3) \to (3, 2) \to (2, 1) \to (1, 0)$.
\end{example}

\begin{proposition}\label{cryptography:proposition-exponentiation}
For $a \in \ZZ/n$ and $e \geq 1$, $a^e$ can be computed with at most $2\log_2e$ multiplications in $\ZZ/n$, hence in $O(\log e\log^2n)$ bit operations.
\end{proposition}

\begin{proof}
Compute the squares $a, a^2, a^4, \dots, a^{2^m}$ with $2^m \leq e < 2^{m+1}$ ($m$ squarings) and multiply those corresponding to the binary digits $1$ of $e$ (at most $m$ multiplications).
\end{proof}

\begin{proposition}[Karatsuba]\label{cryptography:proposition-karatsuba}
Two $k$-bit integers can be multiplied with $O(k^{\log_23}) = O(k^{1.585})$ bit operations.
\end{proposition}

\begin{proof}
For $k = 2m$ write $x = x_1 2^m + x_0$ and $y = y_12^m + y_0$ with $x_i, y_i < 2^m$. Then
\[
xy = x_1y_1 2^{2m} + \bigl((x_1 + x_0)(y_1 + y_0) - x_1y_1 - x_0y_0\bigr)2^m + x_0y_0,
\]
which needs three products of numbers of about $m$ bits and $O(k)$ additions and shifts. So the cost $T(k)$ satisfies $T(2m) \leq 3T(m + 1) + ck$. For $k = 2^j$ (padding with zeros) unwinding gives
$T(2^j) \leq 3^jT(1) + c^{\prime}\sum_{i<j}3^i2^{j-i} = O(3^j) = O(k^{\log_23})$.
\end{proof}

Using the fast Fourier transform, Schönhage and Strassen multiplied in $O(k\log k\log\log k)$ \cite{crandall-pomerance}, and Harvey and van der Hoeven reached $O(k\log k)$
\cite{harvey-van-der-hoeven-2021}.

\section{Primality}\label{cryptography:section-primality}

Fermat's little theorem says that $a^{n-1} \equiv 1 \pmod n$ for prime $n$ and $\gcd(a, n) = 1$ (Corollary~\ref{groups:corollary-euler-fermat}). Its converse fails for the Carmichael numbers
(Example~\ref{groups:example-euler-fermat}), which pass this test for every base.

\begin{proposition}[Korselt]\label{cryptography:proposition-korselt}
A composite $n$ satisfies $a^{n-1} \equiv 1 \pmod n$ for all $a$ coprime to $n$ (it is a \emph{Carmichael number}) if and only if $n$ is squarefree and $p - 1 \mid n - 1$ for every prime $p \mid n$. A Carmichael number
is odd and has at least three prime factors.
\end{proposition}

\begin{proof}
\emph{If.} For $p \mid n$ and $\gcd(a, n) = 1$, $a^{n-1} = (a^{p-1})^{(n-1)/(p-1)} \equiv 1 \pmod p$, and the distinct primes dividing $n$ multiply to $n$ (Theorem~\ref{number-theory:theorem-crt}).

\emph{Only if.} If $p^2 \mid n$, then $a = 1 + n/p$ satisfies $a^j \equiv 1 + jn/p \pmod n$ (as $n \mid (n/p)^2$), so $a$ has order $p$ in $U_n$; then $p \mid n - 1$, impossible as $p \mid n$. So $n$ is squarefree. For
$p \mid n$, let $g$ be a primitive root modulo $p$ (Theorem~\ref{number-theory:theorem-primitive-root-prime}) and, by the Chinese remainder theorem, $a \equiv g \pmod p$, $a \equiv 1 \pmod{n/p}$; then
$g^{n-1} \equiv a^{n-1} \equiv 1 \pmod p$, so $p - 1 \mid n - 1$.

\emph{Shape.} If $n$ were even, an odd prime $p \mid n$ would give $p - 1$ even dividing the odd $n - 1$. If $n = pq$ with $p < q$, then $q - 1 \mid n - 1 = p(q - 1) + (p - 1)$ gives $q - 1 \mid p - 1 < q - 1$, absurd.
\end{proof}

The Carmichael numbers begin $561 = 3 \cdot 11 \cdot 17$, $1105$, $1729$, and there are infinitely many of them (Alford, Granville and Pomerance, 1994 \cite{alford-granville-pomerance}). The Miller--Rabin
test repairs the Fermat test by looking at square roots of $1$: modulo a prime, the only square roots of $1$ are $\pm1$.

\begin{definition}\label{cryptography:definition-strong-liar}
Let $n$ be odd and $n - 1 = 2^st$ with $t$ odd. An integer $a$ coprime to $n$ is a \emph{strong liar} for $n$ if $a^t \equiv 1 \pmod n$ or $a^{2^it} \equiv -1 \pmod n$ for some $0 \leq i < s$. For prime $n$ every
$a$ is a strong liar: the sequence $a^t, a^{2t}, \dots, a^{2^st} = a^{n-1} = 1$ ends at $1$, and the term before the first $1$, if any, is a square root of $1$ different from $1$, hence $-1$.
\end{definition}

\begin{lemma}\label{cryptography:lemma-square-roots}
Let $n$ be odd and divisible by two distinct primes, let $m = 2^st$ be even ($s \geq 1$, $t$ odd), and let $L_m$ be the set of $a \in U_n$ with $a^t \equiv 1$ or $a^{2^it} \equiv -1$ for some $0 \leq i < s$. Then $L_m$ is contained in
a proper subgroup of $U_n$; in particular $|L_m| \leq \varphi(n)/2$.
\end{lemma}

\begin{proof}
Let $j$ be the largest integer $< s$ such that $b^{2^jt} \equiv -1$ for some $b \in U_n$; it exists, since $(-1)^t = -1$. Let $B = \{x \in U_n : x^{2^jt} \equiv \pm1\}$, a subgroup. \emph{$L_m \subseteq B$:} if
$a^t \equiv 1$, then $a^{2^jt} \equiv 1$; if $a^{2^it} \equiv -1$, then $i \leq j$ by maximality and $a^{2^jt} = (a^{2^it})^{2^{j-i}} \equiv \pm1$. \emph{$B \neq U_n$:} write $n = n_1n_2$ with $n_1$, $n_2 > 1$ coprime
(possible as $n$ has two distinct prime factors), and by the Chinese remainder theorem choose $w \equiv b \pmod{n_1}$, $w \equiv 1 \pmod{n_2}$. Then $w^{2^jt} \equiv -1 \pmod{n_1}$ and $\equiv 1 \pmod{n_2}$, so
$w^{2^jt} \not\equiv \pm1 \pmod n$. A proper subgroup has index at least $2$.
\end{proof}

\begin{theorem}[Miller--Rabin]\label{cryptography:theorem-miller-rabin}
If $n > 2$ is an odd composite number, at most $\varphi(n)/2$ elements of $U_n$ are strong liars for $n$. Hence the test "choose $a$ at random in $[1, n - 1]$ and check whether it is a strong liar", repeated $r$
times independently, declares a composite number prime with probability at most $2^{-r}$.
\end{theorem}

\begin{proof}
If $n$ has two distinct prime factors, apply Lemma~\ref{cryptography:lemma-square-roots} with $m = n - 1$. If $n = p^k$ with $k \geq 2$, every strong liar is a Fermat liar, $a^{n-1} \equiv 1$, and the Fermat
liars form a subgroup of the cyclic group $U_{p^k}$ (Theorem~\ref{number-theory:theorem-primitive-root-prime-power}) of order $\gcd(n - 1, p^{k-1}(p - 1)) = \gcd(n - 1, p - 1) < \varphi(n)$, as $p \nmid n - 1$; so it is a
proper subgroup. A random $a$ not coprime to $n$ reveals compositeness through $\gcd(a, n)$.
\end{proof}

A finer count gives the bound $\varphi(n)/4$ \cite{crandall-pomerance}. Each round costs one exponentiation, so the test runs in polynomial time. If the generalised Riemann hypothesis holds, every odd
composite $n$ has a non-liar $a < 2\log^2n$ (Bach), which would make the test deterministic \cite{crandall-pomerance}.

\begin{theorem}[Lucas, Lehmer]\label{cryptography:theorem-lucas}
An integer $n \geq 2$ is prime if and only if there is $a$ with $a^{n-1} \equiv 1 \pmod n$ and $a^{(n-1)/q} \not\equiv 1 \pmod n$ for every prime $q \mid n - 1$.
\end{theorem}

\begin{proof}
If $n$ is prime, take a primitive root. Conversely, the conditions say that the order of $a$ in $U_n$ divides $n - 1$ but no $(n - 1)/q$, so it equals $n - 1$. Then $\varphi(n) \geq n - 1$, which forces $n$ prime.
\end{proof}

Given the factorisation of $n - 1$ and the base $a$, the conditions are checked quickly; applying the same to the prime factors of $n - 1$ recursively gives a certificate of primality of polynomial size
(Pratt, 1975), so PRIMES lies in $\mathrm{NP}$. Much more is true.

\begin{lemma}\label{cryptography:lemma-aks}
Let $n \geq 2$ and $\gcd(a, n) = 1$. Then $n$ is prime if and only if $(x + a)^n \equiv x^n + a$ in $(\ZZ/n)[x]$.
\end{lemma}

\begin{proof}
The coefficient of $x^i$ in $(x + a)^n - x^n - a$ is $\binom nia^{n-i}$ for $0 < i < n$, and $a^n - a$ for $i = 0$. If $n$ is prime, $n \mid \binom ni$ for $0 < i < n$ and $a^n \equiv a$. If $n$ is composite, let
$q$ be a prime factor with $q^k \,\|\, n$. Then $\binom nq = \frac{n(n-1)\cdots(n - q + 1)}{q!}$ has $q$-adic valuation $k - 1$ (the factors $n - 1, \dots, n - q + 1$ are prime to $q$), so $q^k \nmid \binom nq$, and as
$a$ is prime to $q$, the coefficient of $x^q$ is not $0$ modulo $n$.
\end{proof}

\begin{theorem}[Agrawal, Kayal, Saxena]\label{cryptography:theorem-aks}
Primality can be decided in deterministic polynomial time \cite{agrawal-kayal-saxena}.
\end{theorem}

The algorithm checks the congruence of Lemma~\ref{cryptography:lemma-aks} modulo $x^r - 1$ for a small $r$ (of size $O(\log^5n)$) and for $a$ up to about $\sqrt r\log n$; the proof that this suffices uses only
finite fields and cyclotomic polynomials (Chapter~\ref{fields}).

\section{Factoring}\label{cryptography:section-factoring}

Trial division by all $d \leq \sqrt n$ finds a factor in $O(\sqrt n)$ divisions, exponential in the size of $n$. All faster methods exploit structure.

\begin{lemma}[Congruence of squares]\label{cryptography:lemma-squares}
If $x^2 \equiv y^2 \pmod n$ and $x \not\equiv \pm y \pmod n$, then $\gcd(x - y, n)$ is a nontrivial factor of $n$. If $n$ is odd with at least two distinct prime factors, and $y$ is coprime to $n$, at least half of
the solutions $x$ of $x^2 \equiv y^2$ satisfy $x \not\equiv \pm y$.
\end{lemma}

\begin{proof}
$n$ divides $(x - y)(x + y)$ but neither factor, so it shares a nontrivial factor with each. For the second statement, let $n = \prod_{i=1}^kp_i^{e_i}$ with $k \geq 2$. Modulo each $p_i^{e_i}$ the
congruence $x^2 \equiv y^2$ has exactly the two solutions $x \equiv \pm y$, since $U_{p_i^{e_i}}$ is cyclic of even order (Theorem~\ref{number-theory:theorem-primitive-root-prime-power}) and so has exactly two
square roots of $1$. By the Chinese remainder theorem there are $2^k$ solutions modulo $n$, of which only two are $\pm y$, and $2/2^k \leq \frac12$.
\end{proof}

\begin{example}\label{cryptography:example-fermat-factoring}
\emph{Fermat's method} looks for $n = x^2 - y^2 = (x - y)(x + y)$ by trying $x = \lceil\sqrt n\rceil, \lceil\sqrt n\rceil + 1, \dots$ until $x^2 - n$ is a square; it is fast when $n$ has two factors close to $\sqrt n$. For
$n = 8051$, $x = 90$ gives $8100 - 8051 = 49 = 7^2$, so $8051 = 83 \cdot 97$. \emph{Dixon's method} and the \emph{quadratic sieve} produce congruences of squares systematically: they collect many $x$ for which
$x^2 \bmod n$ factors over a set of small primes, and find a subset whose product is a square by linear algebra over $\mathbb{F}_2$ on the exponent vectors.
\end{example}

\begin{proposition}[Pollard's $p - 1$ method]\label{cryptography:proposition-pollard}
Let $p$ be a prime factor of $n$ such that every prime power dividing $p - 1$ is at most $B$, and let $M = \operatorname{lcm}(1, 2, \dots, B)$. For $a$ coprime to $n$, $p$ divides $\gcd(a^M - 1, n)$.
\end{proposition}

\begin{proof}
$p - 1 \mid M$, so $a^M \equiv 1 \pmod p$ by Fermat's theorem.
\end{proof}

The gcd is a proper factor unless $a^M \equiv 1$ modulo every prime factor of $n$; computing $a^M$ costs $O(B\log B)$ multiplications. The method fails when every prime factor $p$ has $p - 1$ with a large prime
factor; Lenstra's \emph{elliptic curve method} replaces the group $U_p$, of fixed order $p - 1$, by the group of points of a random elliptic curve modulo $p$, whose order varies in $[p + 1 - 2\sqrt p, p + 1 + 2\sqrt p]$
(Chapter~\ref{elliptic-curves-arithmetic}).

\begin{lemma}[Birthday bound]\label{cryptography:lemma-birthday}
If $k$ elements are drawn independently and uniformly from a set with $N$ elements, the probability that they are pairwise distinct is at most $e^{-k(k-1)/2N}$.
\end{lemma}

\begin{proof}
The probability is $\prod_{i=1}^{k-1}(1 - \frac iN) \leq \prod e^{-i/N} = e^{-k(k-1)/2N}$, using $1 - x \leq e^{-x}$.
\end{proof}

\begin{remark}\label{cryptography:remark-pollard-rho}
\emph{Pollard's $\rho$ method} iterates $x \mapsto x^2 + c \pmod n$. Modulo an unknown prime factor $p$ the sequence behaves heuristically like a random sequence, so by the birthday bound it repeats after
about $\sqrt p$ steps; Floyd's trick compares $x_i$ with $x_{2i}$ and computes $\gcd(x_{2i} - x_i, n)$, which reveals $p$ at the first repetition modulo $p$. The expected cost $O(\sqrt p) = O(n^{1/4})$ is a
heuristic, not a theorem. The best known methods are subexponential: with $L_n[\alpha, c] = \exp\bigl((c + o(1))(\log n)^\alpha(\log\log n)^{1-\alpha}\bigr)$, the quadratic sieve runs heuristically in $L_n[\frac12, 1]$
and the number field sieve in $L_n[\frac13, (64/9)^{1/3}]$ \cite{crandall-pomerance}.
\end{remark}

\section{Discrete logarithms}\label{cryptography:section-discrete-log}

\begin{definition}\label{cryptography:definition-discrete-log}
Let $G$ be a finite cyclic group, written multiplicatively, with generator $g$ of order $N$. The \emph{discrete logarithm} of $h \in G$ is the unique $x \in \ZZ/N$ with $g^x = h$.
\end{definition}

\begin{proposition}[Baby-step giant-step]\label{cryptography:proposition-bsgs}
Discrete logarithms in $G$ can be computed with $O(\sqrt N)$ group operations and $O(\sqrt N)$ storage.
\end{proposition}

\begin{proof}
Let $m = \lceil\sqrt N\rceil$. Every $x \in [0, N)$ is $x = im + j$ with $0 \leq i, j < m$. Store the \emph{baby steps} $g^j$, $0 \leq j < m$, and compute the \emph{giant steps} $hg^{-im}$, $i = 0, 1, \dots$; the first
coincidence $hg^{-im} = g^j$ gives $x = im + j$. Sorting or hashing the table makes each lookup cheap.
\end{proof}

\begin{proposition}[Pohlig--Hellman]\label{cryptography:proposition-pohlig-hellman}
If $N = \prod q_i^{e_i}$, a discrete logarithm in $G$ can be computed from $\sum_ie_i$ discrete logarithms in subgroups of prime orders $q_i$, plus a polynomial amount of work. So discrete logarithms are
easy when $N$ has only small prime factors.
\end{proposition}

\begin{proof}
\emph{Reduction to prime powers.} $x \bmod q^e$ is the discrete logarithm of $h^{N/q^e}$ to the base $g^{N/q^e}$, which has order $q^e$; the Chinese remainder theorem recovers $x \bmod N$ from the residues.

\emph{Reduction to prime order.} In a group of order $q^e$ generated by $\gamma$, write $x = x_0 + x_1q + \dots + x_{e-1}q^{e-1}$ with digits $0 \leq x_k < q$. Then $h^{q^{e-1}} = (\gamma^{q^{e-1}})^{x_0}$ determines
$x_0$ as a logarithm in the subgroup of order $q$ generated by $\gamma^{q^{e-1}}$; replacing $h$ by $h\gamma^{-x_0}$ and raising to $q^{e-2}$ determines $x_1$, and so on.
\end{proof}

For $G = \mathbb{F}_p^\times$ the \emph{index calculus} method, analogous to the quadratic sieve, is subexponential; for the groups of points of suitable elliptic curves no subexponential method is known,
which is why elliptic curve cryptography uses much shorter keys.

\section{Public-key cryptography}\label{cryptography:section-public-key}

\begin{definition}\label{cryptography:definition-diffie-hellman}
In the \emph{Diffie--Hellman key exchange} \cite{diffie-hellman-1976}, a cyclic group $G = \langle g\rangle$ of order $N$ is public. Alice chooses a secret $a \in \ZZ/N$ and publishes $g^a$; Bob chooses a secret $b$ and
publishes $g^b$. Both compute the shared key $g^{ab} = (g^b)^a = (g^a)^b$. In \emph{ElGamal encryption}, Alice's public key is $h = g^a$; to send $m \in G$, Bob picks a random $k$ and sends $(g^k, mh^k)$, and
Alice recovers $m = mh^k \cdot (g^k)^{-a}$.
\end{definition}

An eavesdropper sees $g$, $g^a$, $g^b$ and must compute $g^{ab}$ (the \emph{computational Diffie--Hellman problem}), which is no harder than computing discrete logarithms; by
Proposition~\ref{cryptography:proposition-pohlig-hellman}, $N$ must have a large prime factor.

For RSA (Proposition~\ref{number-theory:proposition-rsa}, \cite{rivest-shamir-adleman}), with public key $(n, e)$ and secret $d$, finding $d$ is as hard as factoring $n$: knowing $\varphi(n)$ gives $p$ and $q$
(Exercise~\ref{number-theory:exercise-rsa-factoring}), and even knowing $d$ suffices.

\begin{proposition}\label{cryptography:proposition-rsa-factoring}
Given $n = pq$ (distinct odd primes) and $e$, $d$ with $ed \equiv 1 \pmod{\varphi(n)}$, a randomised algorithm factors $n$ in expected polynomial time.
\end{proposition}

\begin{proof}
Let $m = ed - 1 = 2^st$ with $t$ odd; $\varphi(n) \mid m$, so $a^m \equiv 1$ for all $a \in U_n$. Choose $a$ at random; if $\gcd(a, n) > 1$ we are done. If $a \notin L_m$ (notation of
Lemma~\ref{cryptography:lemma-square-roots}), the sequence $a^t, a^{2t}, \dots, a^{2^st} = 1$ does not start with $1$, and the term $x$ before its first $1$ satisfies $x^2 \equiv 1$, $x \not\equiv \pm1$; then
$\gcd(x - 1, n)$ is a factor by Lemma~\ref{cryptography:lemma-squares}. By Lemma~\ref{cryptography:lemma-square-roots} this happens with probability at least $\frac12$ for each $a$.
\end{proof}

\begin{example}[Håstad's broadcast attack]\label{cryptography:example-hastad}
Suppose the same message $m$ is sent with exponent $e = 3$ to three users with pairwise coprime moduli $n_1$, $n_2$, $n_3$, as $c_i \equiv m^3 \pmod{n_i}$. The Chinese remainder theorem gives
$C \equiv m^3 \pmod{n_1n_2n_3}$ with $0 \leq C < n_1n_2n_3$; since $m < n_i$ for each $i$, $m^3 < n_1n_2n_3$, so $C = m^3$ as integers and $m = \sqrt[3]{C}$ is found without factoring. This is why messages are
padded with random bits before encryption; Coppersmith's method, based on the LLL algorithm (Section~\ref{geometry-of-numbers:subsection-lll}), extends such attacks to small roots of polynomial congruences.
\end{example}

\begin{remark}\label{cryptography:remark-signatures-quantum}
\emph{Signatures}: with RSA, the owner of $d$ signs a (hashed) message $m$ as $s = m^d \bmod n$, and anyone verifies $s^e \equiv m$. \emph{Quantum computers}: Shor's algorithm factors integers and computes
discrete logarithms in polynomial time on a quantum computer, by finding periods of $x \mapsto a^x \bmod n$ with the quantum Fourier transform \cite{shor-1997}. This motivates \emph{post-quantum} systems,
such as those based on the \emph{learning with errors} problem, whose hardness is related to worst-case problems on lattices \cite{regev-2009}.
\end{remark}

\section{Exercises}\label{cryptography:section-exercises}

\begin{exercise}\label{cryptography:exercise-carmichael}
Show with Korselt's criterion that $1105 = 5 \cdot 13 \cdot 17$ and $1729 = 7 \cdot 13 \cdot 19$ are Carmichael numbers.
\end{exercise}

\begin{solution}
$1104 = 2^4 \cdot 3 \cdot 23$ is divisible by $4$, $12$ and $16$; $1728 = 2^6 \cdot 3^3$ is divisible by $6$, $12$ and $18$. Both numbers are squarefree, so Proposition~\ref{cryptography:proposition-korselt} applies.
\end{solution}

\begin{exercise}\label{cryptography:exercise-strong-liar-561}
Show that $2$ is not a strong liar for $561$, although $2^{560} \equiv 1 \pmod{561}$.
\end{exercise}

\begin{solution}
$560 = 2^4 \cdot 35$. Modulo $3$, $11$, $17$: $2^{35} \equiv 2, 10, 8$ (as $2^2 \equiv 1 \pmod 3$, $2^5 \equiv -1 \pmod{11}$, $2^8 \equiv 1 \pmod{17}$ and $2^3 = 8$). So $2^{35} \equiv 263 \pmod{561}$, which is not $\pm1$.
Squaring: $2^{70} \equiv 1, 1, 13$, $2^{140} \equiv 1, 1, 16$, $2^{280} \equiv 1, 1, 1$ modulo $3, 11, 17$. So $2^{140} \equiv 67 \pmod{561}$ is a square root of $1$ different from $\pm1$, and $\gcd(67 - 1, 561) = 33$ is a factor.
\end{solution}

\begin{exercise}\label{cryptography:exercise-pollard}
Show, without computing $2^{60} \bmod 1403$, that Pollard's $p - 1$ method with $B = 5$ and $a = 2$ factors $1403 = 23 \cdot 61$.
\end{exercise}

\begin{solution}
$M = \operatorname{lcm}(1, \dots, 5) = 60$. As $60 = 61 - 1$, $2^{60} \equiv 1 \pmod{61}$. Modulo $23$, the order of $2$ is $11$ ($2^{11} = 2048 = 89 \cdot 23 + 1$), and $11 \nmid 60$, so $2^{60} \not\equiv 1 \pmod{23}$. Hence
$\gcd(2^{60} - 1, 1403) = 61$.
\end{solution}

\begin{exercise}\label{cryptography:exercise-bsgs}
Solve $2^x \equiv 13 \pmod{23}$ by baby-step giant-step, using that $2$ has order $11$ modulo $23$.
\end{exercise}

\begin{solution}
$m = 4$. Baby steps $2^j$, $j = 0, \dots, 3$: $1, 2, 4, 8$. Giant steps use $2^{-4} \equiv 13$ (as $16 \cdot 13 = 208 = 9 \cdot 23 + 1$): $13$, $13 \cdot 13 = 169 \equiv 8$. The coincidence $13 \cdot 2^{-4} \equiv 8 = 2^3$
gives $x = 4 + 3 = 7$; indeed $2^7 = 128 = 5 \cdot 23 + 13$.
\end{solution}

\begin{exercise}\label{cryptography:exercise-common-modulus}
Two users share an RSA modulus $n$ with coprime public exponents $e_1$, $e_2$, and the same message $m$ is sent to both. Show how an eavesdropper recovers $m$.
\end{exercise}

\begin{solution}
With Bézout, $ue_1 + ve_2 = 1$ (one of $u$, $v$ negative; invert the corresponding ciphertext modulo $n$, which is possible unless it reveals a factor of $n$). Then
$c_1^uc_2^v \equiv m^{ue_1 + ve_2} = m \pmod n$.
\end{solution}
